% Copyright (c) 2026 Eugene Joseph M. Ragasa.
% SPDX-License-Identifier: Apache-2.0
% Author: Eugene Joseph M. Ragasa
% Formal proof completed: 2026-10-20
\documentclass[11pt]{article}
\usepackage[T1]{fontenc}\usepackage{amsmath,amssymb}\usepackage[margin=1in]{geometry}\usepackage[hidelinks]{hyperref}
\title{Proof Dossier: Global-Phase Invariance of Finite Born Weights}\author{Eugene Joseph M. Ragasa}\date{Formal proof completed: 2026-10-20}
\begin{document}\maketitle
\section{Identity and status}
ID \texttt{QM-MEASUREMENT-THM-002}; status \texttt{proved}. Lean declaration:
\par\noindent{\scriptsize\path{ProjectKoios.QuantumMechanics.Measurement.bornProbability_phaseShift}}.
\section{Exact statement}
For a bounded complex-linear outcome operation $M_i$, normalized representative $\psi$, and
complex scalar $z$ with $\lVert z\rVert=1$,
\[p(i\mid z\psi)=p(i\mid\psi).\]
\section{Source and assumptions}
This is a source-neutral result for the repository's finite bounded-operation definition. An exact
textbook crosswalk is deferred. It proves invariance of the defined probabilities, not equality of
vectors or construction of a quotient space of rays.
\section{Proof schematic and proof}
Complex linearity gives $M_i(z\psi)=zM_i\psi$. Norm homogeneity gives
$\lVert zM_i\psi\rVert=\lVert z\rVert\lVert M_i\psi\rVert$. The unit-modulus premise removes the
first factor, and squaring preserves the equality.
\section{Lean realization and trusted base}
The Lean proof unfolds \texttt{bornProbability} and \texttt{phaseShift}, then simplifies with
bounded-map scalar linearity and \texttt{norm\_smul}. Kernel-reported axioms are
\texttt{propext}, \texttt{Classical.choice}, and \texttt{Quot.sound}; there are no admissions.
\section{Boundary and review}
No physical ray quotient, state-preparation equivalence, experiment, or numerical result is
established. AI-assisted and kernel-checked; no independent expert review is recorded.
\end{document}
